\chapter{Purpose}

\eyebrow{why this software exists}

\vspace{4pt}
Apex is a desktop application for the \textbf{collection and analysis of ovine
carcass images}, built to assemble a research dataset and continue the carcass
typification research at iCEV.

Its purpose is narrow and deliberate. A previous data-collection effort did not
fail because of its machine-learning methods --- it failed because of the
\textbf{integrity of the pairing} between each image, the animal it came from, and
the grade assigned to it. That link was not recorded at the moment of collection;
it was reconstructed afterwards, by file order, without verification. Physical
tags were duplicated or illegible, the laboratory spreadsheet had no image-name
column, and only a few tags matched the lab records. The result was a dataset in
which no image could be trusted to belong to the animal it was attributed to.

\begin{note}[title=The core idea]
Apex is not merely an image capturer. It is the tool that makes the
image\,$\leftrightarrow$\,animal\,$\leftrightarrow$\,grade pairing \textbf{verifiable
by construction}, so that the earlier failure cannot be repeated. Every
architectural decision serves that goal.
\end{note}

\begin{note}[title={A note on the name}]
The application is named \textbf{Apex}. Its window still carries the internal mark
\tele{CARCASS} in the title bar and status bar (visible in the screenshots), which
is the instrument's on-screen signature; throughout this manual the product is
referred to as Apex.
\end{note}

\section{What it does}

Around that integrity core, the software provides an end-to-end workflow:

\begin{itemize}[leftmargin=1.3em,itemsep=3pt]
  \item \textbf{Acquisition.} Register a carcass with a mandatory physical tag,
        then capture images from an external camera --- the image is born paired
        to the carcass in the database, never inferred from a filename.
  \item \textbf{Import.} Bring in photos captured elsewhere; each becomes a
        carcass (or is reconciled to an existing one), so external images enter
        the dataset without breaking pairing.
  \item \textbf{Analysis.} Run trained models over the images: fat segmentation
        (validated), an experimental finishing grade, and an analytical
        conformation (convexity) measure --- all with automatic background
        removal so surface measurements are correct.
  \item \textbf{Grading.} Collect independent, blind grades from multiple raters
        and measure their agreement (Fleiss' $\kappa$).
  \item \textbf{Live monitor.} A real-time fat overlay from a camera or a video
        file, to help frame the carcass during capture.
  \item \textbf{Dashboard \& export.} Track sample progress against the target,
        and export the dataset as a manifest plus an integrity report.
\end{itemize}

\section{Who it is for}

The primary users are the research team and the operators who collect carcass
images at the slaughter line and in the laboratory. The interface is designed as
a scientific \emph{instrument} --- scanned and operated, not read top to bottom
--- and follows the visual language of mission-control software: a dark
instrument ground, a cyan signal accent, and monospace readouts for every
measured value.

\section{A note on scientific honesty}

Two of the numbers this software produces --- the \emph{finishing grade} and the
\emph{conformation grade} --- are marked throughout the interface as
\textbf{experimental / not validated}. This is a deliberate choice, grounded in
the project's own findings (Station~13). The software surfaces them because they
are useful as exploratory references and because the dataset it builds is exactly
what is needed to validate them in the future. It never presents them as settled
facts.
